\section{Cronograma}

\textbf{Extensión máxima: 2 páginas}

Presente un cronograma realista para el desarrollo del proyecto en 5 meses (192
horas). Esto implica una dedicación aproximada de 9 horas a la semana.

\begin{itemize}
  \item Divida el trabajo en fases o entregables mensuales.

  \item Incluya actividades clave alineadas con los objetivos.

  \item Puede usarse una tabla o diagrama de Gantt.
\end{itemize}

Le recomendamos seguir esta guía:

\begin{enumerate}
  \item \textbf{Seleccione y adapte una metodología para el desarrollo del
    proyecto.Ejemplo}
    \begin{itemize}
      \item \emph{En cascada}: fases secuenciales (requisitos → diseño →
        ejecución → verificación → cierre).

      \item \emph{Iterativa}: ciclos cortos que repasan análisis, diseño y
        pruebas varias veces.

      \item \emph{Ágil/Scrum}: divida en sprints de 2–4 semanas con entregables incrementales.
    \end{itemize}

  \item \textbf{Defina fases o sprints.} A cada fase o sprint asígnese un
    objetivo y un entregable claro.

  \item \textbf{Estime horas por fase.} Reparta las 192 h entre fases en
    proporción al esfuerzo previsto. P. ej.:
    \begin{itemize}
      \item Investigación y planificación: 10 \% (aprox 19 h)

      \item Ejecución principal (diseño/prototipo/experimento): 50 \% (aprox 96 h)

      \item Evaluación y ajustes: 15 \% (aprox 29 h)

      \item Reuniones con el director, documentación y defensa: 25 \% (aprox 48
        h)
    \end{itemize}

  \item \textbf{Distribuya semanalmente.} Con base en la estimación por fase, calcule
    horas/semana y asigne las tareas de 4–10 h cada una

  \item \textbf{Desglose en tareas concretas.} Para cada fase/sprint, liste actividades
    (p. ej.: lectura de artículos, diseño de experimento, programación de prototipo,
    análisis de datos, redacción de capítulos) y señale:
    \begin{itemize}
      \item Fecha de inicio

      \item Duración estimada (h)
    \end{itemize}
    %\item \textbf{Incluya holgura (\≈10 \%).}
    Reserve ~20 h para imprevistos o refinamientos.

  \item \textbf{Establezca controles periódicos.} Agende revisiones cada 4–6
    semanas (o al cierre de cada sprint) para:
    \begin{itemize}
      \item Comparar horas planificadas vs. reales

      \item Ajustar el cronograma

      \item Recoger feedback de director - codirector
    \end{itemize}

  \item \textbf{Visualice y monitorice.} Use diagrama de Gantt, tablero Kanban o
    tabla de seguimiento en hoja de cálculo.
\end{enumerate}

\medskip
\noindent
\textbf{Ejemplo de desglose semanal (primer mes, 9 h/semana):}
\begin{table}[h!]
  \centering
  \begin{tabular}{p{2cm} p{4cm} p{2cm} p{2cm}}
    \toprule \textbf{Semana} & \textbf{Actividad}                                 & \textbf{Horas} & \textbf{Saldo Mes} \\
    \midrule 1               & Revisión bibliográfica y marco teórico             & 9              & 29                 \\
    2                        & Diseño de cuestionario / experimento inicial       & 9              & 20                 \\
    3                        & Revisión seguimiendo y metodología con el director & 8              & 12                 \\
    4                        & Ajustes y planificación detallada                  & 8              & 4                  \\
    \bottomrule
  \end{tabular}
  \caption{Desglose de 34 h del Mes 1 (incluye 10 \% holgura)}
\end{table}

\bigskip
\textbf{Consejos finales:}
\begin{itemize}
  \item \emph{Transparencia}: Documente cambios (qué, por qué, cuándo). Esto le
    será muy útil para escribir su documento final.

  \item \emph{Flexibilidad}: Recalibre fases según el progreso real.

  \item \emph{Comunicación}: Comparta avances regulares con el director,
    normalmente los avances semanales permiten mejor avance en los proyectos. Si
    tiene algun problema de comunicación con el director durante la ejecución del
    proyecto, repórtelo a la dirección del programa para buscar en conjunto
    soluciones.
\end{itemize}